\chapter{Complexity — U Path Loops}
%%% generated by the blueprint pipeline from TCSlib.Complexity.Randomized.UPathLoops
\section{Overview}

Declarations of the Lean module \texttt{TCSlib.Complexity.Randomized.UPathLoops}.

\section{Declarations}

\begin{theorem}[Single addition in a probabilistic program]
\lean{Complexity.LogProg.psem_addOnce}
\label{Complexity.LogProg.psem_addOnce}
\leanok
\uses{Complexity.LogProg.PSem, Complexity.LogProg.PSem.inc, Complexity.LogProg.PSem.jeq_neg, Complexity.LogProg.PSem.jeq_pos, Complexity.LogProg.RegLe, Complexity.LogProg.register_update_le}
\difficulty{6}
Let $A$ be a nondeterministic register program with $m$ natural-number registers, state
set $\Lambda$, and oracle, and fix an input $x$ and register bound $B$. Let $j,y,R$ be
register indices with $j\ne y$, $j\ne R$, and $R\ne y$. Suppose the instruction at state
$test$ compares registers $j$ and $y$, going to $ex$ when they are equal and to $i_1$
otherwise; at $i_1$ it increments register $R$ and goes to $i_2$, and at $i_2$ it
increments register $j$ and returns to $test$. For any rational $q$ and register
valuation $v$ bounded by $B$, assume $v(j)\le v(y)$ and $v(R)+(v(y)-v(j))\le B$. If the
configuration at $ex$ with register $R$ increased by $v(y)-v(j)$ and register $j$ set to
$v(y)$ has probabilistic semantics of acceptance probability $q$ under the register
bound $B$, then the configuration at $test$ with valuation $v$ has the same semantics
and acceptance probability.
\end{theorem}

\begin{theorem}[Multiplication in a probabilistic program]
\lean{Complexity.LogProg.psem_mul}
\label{Complexity.LogProg.psem_mul}
\leanok
\uses{Complexity.LogProg.PSem, Complexity.LogProg.PSem.clr, Complexity.LogProg.PSem.inc, Complexity.LogProg.PSem.jeq_neg, Complexity.LogProg.PSem.jeq_pos, Complexity.LogProg.RegLe, Complexity.LogProg.psem_addOnce, Complexity.LogProg.register_update_le}
\difficulty{6}
Let $A$ be a nondeterministic register program with $m$ natural-number registers, state
set $\Lambda$, and oracle, and fix an input $x$ and register bound $B$. Let $i,j,a,b,R$
be register indices satisfying $i\ne a$, $j\ne b$, $j\ne R$, $R\ne b$, $i\ne j$, $i\ne
R$, $R\ne a$, $j\ne a$, and $i\ne b$. Suppose the instructions are as follows: $mo$
compares $i$ and $a$, going to $ex$ if equal and to $mo_1$ otherwise; $mo_1$ clears
register $j$ and goes to $mi$; $mi$ compares $j$ and $b$, going to $mo_2$ if equal and
to $mi_1$ otherwise; $mi_1$ increments $R$ and goes to $mi_2$; $mi_2$ increments $j$ and
returns to $mi$; and $mo_2$ increments $i$ and returns to $mo$. For any rational $q$ and
valuation $v$ bounded by $B$ with $v(i)\le v(a)$ and $v(R)+(v(a)-v(i))v(b)\le B$, assume
that for every $z\le B$, the configuration at $ex$ with $R$ set to
$v(R)+(v(a)-v(i))v(b)$, $i$ set to $v(a)$, and $j$ set to $z$ has probabilistic
semantics with acceptance probability $q$. Then the configuration at $mo$ with valuation
$v$ has the same semantics and acceptance probability.
\end{theorem}

\begin{lemma}[Copy-site semantics preservation]
\lean{Complexity.LogProg.ups_cp}
\label{Complexity.LogProg.ups_cp}
\leanok
\uses{Complexity.LogProg.CpSite, Complexity.LogProg.PSem, Complexity.LogProg.RegLe, Complexity.LogProg.ULabel, Complexity.LogProg.UPS, Complexity.LogProg.UReg, Complexity.LogProg.UReg.Le, Complexity.LogProg.UReg.reg, Complexity.LogProg.UReg.v, Complexity.LogProg.cpDst, Complexity.LogProg.cpEx, Complexity.LogProg.cpSrc, Complexity.LogProg.pathOr, Complexity.LogProg.psem_copy, Complexity.LogProg.upA}
\difficulty{1}
Let $x$ be a Boolean input, $B$ a natural-number register bound, $k$ a copy site, $p$ a
rational number, and $\sigma$ a register state. Suppose every register of $\sigma$ is at
most $B$, and the value in the destination register selected by $k$ is at most the value
in its source register. If the program has probabilistic semantics $p$ from the entry
state associated with $k$, with the destination register replaced by the source
register’s value, then it has probabilistic semantics $p$ from the copy instruction for
$k$ with the original register state.
\end{lemma}

\begin{lemma}[Address Column Distinctness]
\lean{Complexity.LogProg.adSite_ne}
\label{Complexity.LogProg.adSite_ne}
\leanok
\uses{Complexity.LogProg.AdSite, Complexity.LogProg.adP, Complexity.LogProg.adY}
\difficulty{4}
For every address site $k$, the column indexed by $3$ in the set of $15$ columns differs
from each of the columns assigned to $k$ by the two address maps, and those two assigned
columns differ from each other.
\end{lemma}

\begin{lemma}[Arithmetic-site semantics preservation]
\lean{Complexity.LogProg.ups_ad}
\label{Complexity.LogProg.ups_ad}
\leanok
\uses{Complexity.LogProg.AdSite, Complexity.LogProg.PSem, Complexity.LogProg.RegLe, Complexity.LogProg.ULabel, Complexity.LogProg.UPS, Complexity.LogProg.UReg, Complexity.LogProg.UReg.Le, Complexity.LogProg.UReg.reg, Complexity.LogProg.UReg.v, Complexity.LogProg.adEx, Complexity.LogProg.adP, Complexity.LogProg.adSite_ne, Complexity.LogProg.adY, Complexity.LogProg.pathOr, Complexity.LogProg.psem_addTwice, Complexity.LogProg.upA}
\difficulty{1}
Let $x$ be a Boolean word, $B$ a natural-number register bound, $k$ an arithmetic site,
$q$ a rational number, and $\sigma$ a register state. Suppose every register of $\sigma$
is at most $B$, its inner counter $b$ is at most the value of the register selected by
$k$, and the value of the register selected as the scratch destination plus twice the
difference between that selected comparison value and $b$ is at most $B$. If the program
has probabilistic semantics $q$ from the site’s associated entry label, with the scratch
destination first set to that scratch value plus twice the difference and register $3$
then set to the comparison value, then it has probabilistic semantics $q$ from the
arithmetic instruction for $k$ with the original register state.
\end{lemma}

\begin{lemma}[Multiplication subroutine semantics]
\lean{Complexity.LogProg.ups_mul}
\label{Complexity.LogProg.ups_mul}
\leanok
\uses{Complexity.LogProg.MulSite, Complexity.LogProg.PSem, Complexity.LogProg.PSem.clr, Complexity.LogProg.RegLe, Complexity.LogProg.ULabel, Complexity.LogProg.UPS, Complexity.LogProg.UReg, Complexity.LogProg.UReg.Le, Complexity.LogProg.UReg.reg, Complexity.LogProg.UReg.v, Complexity.LogProg.mulA, Complexity.LogProg.mulB, Complexity.LogProg.mulEx, Complexity.LogProg.mulR, Complexity.LogProg.pathOr, Complexity.LogProg.psem_mul, Complexity.LogProg.register_update_le, Complexity.LogProg.upA}
\difficulty{5}
Let $x$ be a finite bit string, $B$ a natural-number bound, $k$ one of the two
multiplication sites, $q$ a rational number, and $\sigma$ a register state. Suppose
every register of $\sigma$ is at most $B$, and the value in the result register plus the
product of the two input-register values is at most $B$. Suppose also that for every
natural number $z\le B$, the universal probabilistic program started at the exit site
associated with $k$, with its outer counter set to the first input value, its inner
counter set to $z$, and its result register set to the original result value plus the
product of the input values, has probabilistic semantics with acceptance probability
$q$. Then the program started at the entry site associated with $k$ and register state
$\sigma$ has probabilistic semantics with acceptance probability $q$.
\end{lemma}

\begin{lemma}[Excluded row indices]
\lean{Complexity.LogProg.ixRow_ne}
\label{Complexity.LogProg.ixRow_ne}
\leanok
\uses{Complexity.LogProg.IxSite, Complexity.LogProg.ixRow}
\difficulty{4}
For each of the four sites $y_0,y_1,y_2,w$, its assigned row index is different from
$0$, from $2$, and from $3$.
\end{lemma}

\begin{lemma}[Excluded column indices]
\lean{Complexity.LogProg.ixCol_ne}
\label{Complexity.LogProg.ixCol_ne}
\leanok
\uses{Complexity.LogProg.IxSite, Complexity.LogProg.ixCol}
\difficulty{4}
For each of the four sites $y_0,y_1,y_2,w$, its assigned column index is different from
$0$, from $2$, and from $3$.
\end{lemma}

\begin{lemma}[Unaffected row and column values]
\lean{Complexity.LogProg.ix_v_upd}
\label{Complexity.LogProg.ix_v_upd}
\leanok
\uses{Complexity.LogProg.IxSite, Complexity.LogProg.UReg, Complexity.LogProg.UReg.v, Complexity.LogProg.ixCol, Complexity.LogProg.ixRow}
\difficulty{4}
Let $k$ be one of the four index sites, let $\sigma$ be a register state, and let
$p',a',b'$ be natural numbers. After replacing the input pointer and the two counters
$a,b$ of $\sigma$ by $p',a',b'$, respectively, the values of the resulting register
state at the row and column indices assigned to $k$ equal the corresponding values of
$\sigma$.
\end{lemma}

\begin{lemma}[Row Scan Invariant]
\lean{Complexity.LogProg.ups_row}
\label{Complexity.LogProg.ups_row}
\leanok
\uses{Complexity.LogProg.IxSite, Complexity.LogProg.PSem.clr, Complexity.LogProg.PSem.inc, Complexity.LogProg.PSem.jeq_neg, Complexity.LogProg.PSem.jeq_pos, Complexity.LogProg.UPS, Complexity.LogProg.UPS.of_eq, Complexity.LogProg.UReg, Complexity.LogProg.UReg.Le, Complexity.LogProg.UReg.reg, Complexity.LogProg.UReg.v, Complexity.LogProg.adEx, Complexity.LogProg.adP, Complexity.LogProg.adY, Complexity.LogProg.ixRow, Complexity.LogProg.ix_v_upd, Complexity.LogProg.upA, Complexity.LogProg.ups_ad}
\difficulty{6}
Let $x$ be a Boolean string, $B$ a natural-number register bound, $k$ an index site, and
$q$ a rational number. For every natural number $dd$ and register state $\sigma$,
suppose the value in the row register selected by $k$ minus the outer counter, with
truncated natural subtraction, equals $dd$; every register is at most $B$; the outer
counter is at most the selected row-register value, which is less than the vertex count
$n$; the input pointer is $2n+2+2(an)$, where $a$ is the outer counter; the inner
counter is at most $n$; and $2n+2+2n^2\le B$. If, for every $b'\le n$, the configuration
at the next-row state with outer counter set to the selected row-register value, input
pointer set to $2n+2+2(n\cdot\mathrm{row})$, and inner counter set to $b'$ has
probabilistic semantics with acceptance probability $q$, then the configuration at the
row-check state with register state $\sigma$ has probabilistic semantics with acceptance
probability $q$.
\end{lemma}

\begin{lemma}[Index routine preserves probabilistic semantics]
\lean{Complexity.LogProg.ups_index}
\label{Complexity.LogProg.ups_index}
\leanok
\uses{Complexity.LogProg.IxSite, Complexity.LogProg.PSem.clr, Complexity.LogProg.PSem.inc, Complexity.LogProg.UPS, Complexity.LogProg.UPS.of_eq, Complexity.LogProg.UReg, Complexity.LogProg.UReg.Le, Complexity.LogProg.UReg.reg, Complexity.LogProg.UReg.v, Complexity.LogProg.UReg.v_0, Complexity.LogProg.adEx, Complexity.LogProg.adP, Complexity.LogProg.adY, Complexity.LogProg.ixCol, Complexity.LogProg.ixEx, Complexity.LogProg.ixRow, Complexity.LogProg.ix_v_upd, Complexity.LogProg.upA, Complexity.LogProg.ups_ad, Complexity.LogProg.ups_row, Complexity.rowMajor_lt}
\difficulty{6}
Let $x$ be a bit string, $B$ a natural number, $q$ a rational number, $k$ one of the
four index sites, and $\sigma$ a register state whose fifteen values are at most $B$.
Write $r$ and $c$ for the values in the row and column positions specified by $k$, and
suppose $r<n$ and $c<n$, where $n$ is the vertex-count register of $\sigma$. Assume
$2n+2+2n^2\le B$, and that the program started at the entry specified by $k$ from the
state obtained by setting its input pointer to $2n+2+2(rn+c)$ and its row and inner
counters to $r,c$ has all live register values at most $B$, terminates along every
sequence of choices, and has acceptance probability $q$ and rejection probability $1-q$
for every choice length at which every sequence of choices halts. Then the program
started at the index routine for $k$ from $\sigma$ has those same bounded-register
termination and probability properties.
\end{lemma}
